% design goals
% - "showing the seams", see the internal representation
% - more immediacy, direct feeback
% - feels like an instrument, could be used for composition / performance
% - support itareative designing process
% - see many possibilities at once (casual creators, mutant shopping)
% - hand authoring and tool based generation

% why, what, how, evalution
% answer RQ2: how much can be done in Dusa vs JS
% distribution of code, decoupling UI from solver
% dev process, improving the prototype
% effort between frontend and backend, join points with dusa

\section{UI Implementation}

With the goal in mind of creating a music-specific interface for the Dusa
counterpoint generator, we developed a web-based graphical user interface.
The following subsections describe our design goals, the
status of the current prototype from a user perspective,
and the software architecture and development approach.

\subsection{Design Goals}

We wanted to allow, at minimum, the input of notes in a familiar and
ergonomic notation, and the playback of generated voices as actual sound.
We wanted to create a tight feedback loop in the 
sense of ``direct manipulation'' (TODO cite)
so that users could experience using our solver as
part of a creative process without moving manually
between distinct pieces of software, making
incremental changes and seeing their results with
relative immediacy.
We also wanted to fit some of the criteria for
{\em casual creators}~\cite{compton-casual},
specifically accessibility to users with little
prior background in music, and the ability
to support autotelic creativity (for its own
sake rather than towards a productive end).
Finally, we wanted the ability to share our work 
with a wider audience, motivating development for
the web browser. 

There were two key areas in which we needed to make design decisions:
\begin{enumerate}
\item {\bf Musical interface}: independent of our use
of Dusa, there are many possible designs for
authoring, editing, and playback of musical compositions.
\item {\bf Solver interface}: independent of music as an
application domain, there are many possible designs
for calling a solver, displaying its results, and
letting the user navigate them and integrate them
into their composition.
\end{enumerate}

For the musical interface, we decides to adopt a 
hybrid instrument-composition
tool approach, which could be incorporated into a musician's process at the
very early stages of working out a tune, similar to how
one might ``sketch'' a melody on a piano when first conceptualizing
a composition that will ultimately be performed on different instruments.
We do not intend to support music
{\em production}, nor do we currently imagine using this
tool to support live performance (though we have some ideas
for how it could be extended to do so).
Our interface takes inspiration from
Beepbox (cite/url TODO), having just one view with
the piano roll as its most prominent element.

% There are many possible designs for such a system
% depending on the concept for user experience: 
% we might draw inspiration from digital audio 
% workstations (aimed at production), 
% classical composition software, or livecoding
% environments like Haskore (TODO cite), geared
% more towards performance. Ultimately we landed
% on wanting a piano roll-like interface to support
% sequencing and tracking for a fixed-length
% composition, inspired by Beepbox (TODO cite/url),

For the solver interface, we wanted to support
the user {\em iteratively refining} their work, 
combining hand-authored creative decisions with
generated solutions.
We wanted the user to be able to see many 
generated solutions at once to help them understand
the range of choices available.
We wanted the principles of counterpoint, and their
corresponding constraints on the generated results,
to be possible to inspect and toggle.
Finally, we didn't want to just seal Dusa away
into a black-box abstraction; we wanted to allow
the user to ``see the seams'' of the technical
system under the hood, e.g. by inspecting 
the logical representation of counterpoint rules
and the raw predicate output of the solver.

% In particular, we drew inspiration from the
% {\em casual creators} literature, which identifies
% ``mutant shopping'' as a pattern for successfully
% joyful creative tools~\cite{compton-casual}: 

% We want to design a UI that feels to the user like playing an instrument, and
% allows use for composing and performing musical works. We want to show immediate
% feedback to the user in order to support iteratively refining their work, with
% a combination of handmade and generated edits. Another goal is to show many
% possibilites at once, to help the user see what choices they can make, as a form
% of mutant shopping \cite{compton-casual}.

\subsection{Current Prototype}

A screenshot of the CounterChoice prototype, from the user's
perspective, is in figure TODO. It consists of the following
elements:

\paragraph{Notes Editor}
The notes editor is a fully controlled input component that renders a HTML
\verb|<canvas>| element. The interface resembles a piano roll with 14 notes and
8 time steps. Clicking on a note will set it as the cantus firmus for that
time step. Shift-clicking will set the counterpoint at that time.

% (i.e., require
% that solutions contain \verb|cpNote| facts as specified).

\AP{This section is pretty bare, but I feel that Jim might have more to say
    about the implementation of the editor, as he mainly worked on it.}

\paragraph{Forbidden Scenarios}
TODO

\paragraph{Playback and Solve Buttons}
TODO

\paragraph{Solution Grid}
TODO


\subsection{Architectural Overview}


The UI is implemented as a React web app that calls out to a Dusa solver through
a worker thread. The execution is entirely client-side, but
architecturally we can think of the Dusa worker as the backend and
the rendering code as a frontend, appropriately decoupled.


The frontend records changes made to the inputs, including
the cantus firmus, hand-selected counterpoint notes, and 
which rules to enforce. It  
appends the corresponding text to the Dusa program fed as an input to the
solver.
The main \verb|<App />| component
keeps as state the data that gets passed to the solver.
The Dusa API's \verb|solve| function returns information about the query as
well as streams the solutions through a generator. This represents the only
point at which the web code interfaces with a Dusa program or its API.


% TODO: something about how the toggled check marks and
% hand-chosen facts are represented internally? maybe above.

\paragraph{Dusa Worker}
Since Dusa works synchronously while generating a solution, we need to offload
the work off the main thread so that the UI can remain responsive even while
it is working on generating solutions. To achieve this, we spawn a web worker
that sends a message to the main thread every time a solution gets generated.
These messages are then packaged into an async generator so that we can loop
over the solutions with a \verb|for await ... of| loop, which can be used easily
by JS developers who are used to asynchronous development.

The worker takes as input the cantus firmus and the specified notes of the
counterpoint, as well as the rules that should be enabled. It generates strings
corresponding to the relevant \verb|#demand| and \verb|#forbid| rules, and the
facts for specifying notes. Upon generating a solution, it transforms the raw
facts back into the raw notes for each track. To facilitate users being able to
see into the internals of the program, we also pass along the final generated
Dusa programs, and all facts that were generated as part of the solution:
specifcially flagging facts that denote a solution being an invalid
counterpoint. The raw information is not used, but is simply passed along to the
UI to offer the user a look inside the system.


